% !TeX root = ../../Book3.tex
% 纲要标题：### 1.2 根与约化
% 纲要位置：工作记录/计划纲要.md，第 1920 行。
% 纲要细目（写作范围）：
% * 幂零根和根理想；用有限混合乘积与无限变量例子区分逐元素幂零、统一指数和理想幂零
% * 根理想作为素理想的交；把不含指定元素转为避开其幂组成的乘法集
% * 约化环与零维例子；无限域乘积的论证对全部素理想成立

\section{根与约化}
\label{b3:sec:0102}

给定交换含幺环 \(R\) 的理想 \(I\)，一个元素 \(a\) 属于所有包含 \(I\) 的素理想，当且仅当它在商环 \(R/I\) 中的像是幂零元，也就是某次幂 \(a^n\) 落在 \(I\) 中。这一事实把素理想的包含关系与商环中的幂零现象联系起来。不过，这些素理想并不记录最小的幂次 \(n\)，也不能恢复原来的理想 \(I\)。将 \(I\) 换成它的根，保留的正是这组素理想包含信息。

\subsection{幂零根与根理想}
\label{b3:subsec:0102:01}

\begin{definition}\label{b3:def:radical-reduced}
设 \(R\) 为交换含幺环，\(I\) 为 \(R\) 的理想。定义 \(I\) 的\term[lixiang-de-gen]{根}[radical of an ideal]为
\[
\sqrt I=\{\,a\in R:\text{存在 }n\ge1\text{ 使 }a^n\in I\,\}.
\]
若 \(I=\sqrt I\)，则称 \(I\) 为\term[genlixiang]{根理想}[radical ideal]。\(\sqrt{(0)}\) 称为 \(R\) 的\term[miling-gen]{幂零根}[nilradical]，记为 \(\operatorname{Nil}(R)\)。若 \(\operatorname{Nil}(R)=0\)，则称 \(R\) 为\term[yuehua-huan]{约化环}[reduced ring]。
\end{definition}
\symidx{\(\sqrt I\)：理想的根}
\symidx{\(\operatorname{Nil}(R)\)：环的幂零根}

\(\sqrt I\) 确实是理想。若 \(a^m,b^n\in I\)，展开 \((a+b)^{m+n-1}\)，每项至少含 \(a^m\) 或 \(b^n\)，故该幂属于 \(I\)；负号和环元素倍乘也保持条件。还直接得到
\[
I\subseteq\sqrt I,\qquad \sqrt{\sqrt I}=\sqrt I,
\qquad \sqrt{I\cap J}=\sqrt{IJ}=\sqrt I\cap\sqrt J.
\]
最后一式中，若 \(a^m\in I,a^n\in J\)，则 \(a^{m+n}\in IJ\)；另一方向由 \(IJ\subseteq I\cap J\) 给出。对理想根迭代时，只需把两个存在的正指数相乘，便回到 \(I\)。

\ProseParagraphBreak
商环的幂零根为 \(\sqrt I/I\)，所以 \(R/I\) 约化当且仅当 \(I\) 是根理想。特别地，\(R_{\mathrm{red}}=R/\operatorname{Nil}(R)\) 约化，称为 \(R\) 的约化商。它把全部幂零元素送到零，而不要求剩下的非零元素都可逆。
\symidx{\(R_{\mathrm{red}}\)：环的约化商}

\begin{definition}\label{b3:def:nilpotent-ideal}
设 \(R\) 为交换含幺环，\(J\) 为 \(R\) 的理想。对正整数 \(N\)，\(J^N\) 是全部 \(N\) 个 \(J\) 中元素的乘积所生成的理想。若存在正整数 \(N\) 使 \(J^N=0\)，则称 \(J\) 为\term[mi-ling-lixiang]{幂零理想}[nilpotent ideal]；满足该式的最小正整数称为 \(J\) 的\term[miling-zhishu]{幂零指数}[nilpotency index]。
\end{definition}

理想幂零要求任意 \(N\) 个理想元素的混合乘积都为零，强于只对同一个元素取幂。对理想 \(J\)，下面三个条件逐项变强：
\[
\begin{aligned}
&\forall u\in J\ \exists n(u)\ge1, &&u^{n(u)}=0;\\
&\exists N\ge1\ \forall u\in J, &&u^N=0;\\
&\exists N\ge1, &&J^N=0.
\end{aligned}
\]
第一个条件在英文文献中称为 nil ideal，只要求每个元素各自幂零；第二个条件要求元素幂零有统一上界；第三个条件才是理想幂零。最后一个条件蕴含前两个，反向一般不成立。

\ProseParagraphBreak
例如，设 \(k\) 为域，在商环
\[
R=k[\varepsilon,\delta]/(\varepsilon^2,\delta^2)
\]
中仍用 \(\varepsilon,\delta\) 表示变量的像，令 \(J=(\varepsilon,\delta)\)。这个商环有 \(k\)-基 \(1,\varepsilon,\delta,\varepsilon\delta\)，所以混合乘积 \(\varepsilon\delta\ne0\)。因此
\[
J^2=(\varepsilon\delta)\ne0,\qquad J^3=0,
\]
后一式是因为任意三个生成元的乘积都会重复 \(\varepsilon\) 或 \(\delta\)。若 \(\operatorname{char}k=2\)，每个 \(u\in J\) 甚至都满足 \(u^2=0\)，而 \(J^2\) 仍不为零。这说明同一指数下，单个元素的幂与不同元素的混合乘积不能混为一谈。

\begin{proposition}\label{b3:prop:finite-nil-ideal}
设 \(R\) 为交换含幺环，\(J=(a_1,\ldots,a_m)\) 为 \(R\) 的理想，其中 \(m\ge0\)。若 \(a_i^{n_i}=0\)，且 \(n_i\ge1\)（\(1\le i\le m\)），则
\[
J^N=0,\qquad N=1+\sum_{i=1}^m(n_i-1).
\]
因此 Noether 环的幂零根是幂零理想。
\end{proposition}
\begin{proof}
\(m=0\) 时，\(J=0\)、\(N=1\)，结论成立。设 \(m\ge1\)。\(J^N\) 由全部 \(N\) 个生成元的乘积线性生成。若每个 \(a_i\) 出现次数都小于 \(n_i\)，总个数至多为 \(N-1\)，矛盾。所以每个乘积含某个 \(a_i^{n_i}\)，均为零。Noether 环的幂零根作为理想有限生成，每个生成元又幂零，故适用第一项。
\end{proof}

没有有限生成性时，一个理想即使每个元素都幂零，也未必是幂零理想。取域 \(k\)，考虑
\[
R=k[x_1,x_2,\ldots]/(x_1^2,x_2^2,\ldots).
\]
令 \(J=(\bar x_1,\bar x_2,\ldots)\)。每个 \(u\in J\) 只用有限多个变量，设有 \(r\) 个；任何 \(r+1\) 个正次单项式的乘积都会重复一个变量，所以 \(u^{r+1}=0\)。但 \(\bar x_1\cdots\bar x_N\ne0\) 对任意 \(N\) 成立，因为平方关系生成的理想由含某个平方因子的单项式线性生成，不能包含这个不含平方因子的单项式。又因商 \(R/J=k\) 约化，全部幂零元都在 \(J\) 中。因此 \(J=\operatorname{Nil}(R)\) 而 \(J^N\ne0\) 对全部 \(N\) 成立。

\ProseParagraphBreak
若 \(\operatorname{char}k=2\)，这个例子还有更强的性质：每个 \(u\in J\) 都满足 \(u^2=0\)。因为 \(u\) 是有限个正次单项式的线性组合，特征二时平方把有限和变成各项平方之和，而每个正次单项式的平方都为零。所以元素幂零即使有统一上界，也不能保证理想的某个幂为零。

\subsection{根理想的素理想交表示}
\label{b3:subsec:0102:02}

\begin{theorem}\label{b3:thm:radical-prime-intersection}
设 \(R\) 为交换含幺环。任意理想 \(I\subseteq R\) 满足
\begin{equation}\label{b3:eq:radical-prime-intersection}
\sqrt I=\bigcap_{\mathfrak p\supseteq I,\ \mathfrak p\text{ 素}}\mathfrak p.
\end{equation}
空交取为 \(R\)，因而公式也适用于 \(I=R\)。
\end{theorem}
\begin{proof}
若 \(a^n\in I\subseteq\mathfrak p\)，反复使用素性得到 \(a\in\mathfrak p\)，故左边包含于右边。

反过来，对 \(a\notin\sqrt I\)，我们要找一个包含 \(I\) 却不含 \(a\) 的素理想。对素理想 \(\mathfrak p\)，不含 \(a\) 等价于不含 \(a\) 的任何非负整数次幂：正次幂属于 \(\mathfrak p\) 会由素性推出 \(a\in\mathfrak p\)，而 \(1\) 本就不属于真理想。因此取乘法集 \(S=\{1,a,a^2,\ldots\}\)，便将不含 \(a\) 转化为与 \(S\) 不交。由 \(a\notin\sqrt I\) 可知 \(S\cap I=\varnothing\)：没有正次幂落在 \(I\) 中，且 \(I\ne R\)，所以 \(1\notin I\)。\cref{b3:lem:prime-disjoint-multiplicative}给出包含 \(I\) 并避开 \(S\) 的素理想 \(\mathfrak p\)，特别是 \(a\notin\mathfrak p\)。因此 \(a\) 不在右边的交中，证明反向包含。
\end{proof}

取 \(I=0\)，可知一个元素幂零当且仅当它属于全部素理想。于是映射
\[
R\longrightarrow\prod_{\mathfrak p\text{ 素}}R/\mathfrak p,
\qquad a\longmapsto(a+\mathfrak p)_{\mathfrak p}
\]
的核恰为 \(\operatorname{Nil}(R)\)。约化环因而可以嵌入整环的乘积，但本身未必是整环：\(k\times k\) 已给出最简单的反例。

\ProseParagraphBreak
在一般环上，不能把上述交中的素理想直接换成极大理想。以 \(R=k[t]_{(t)}\) 为例，其唯一极大理想为 \((t)R\)，但 \(R\) 是整环，幂零根为零。因此所有极大理想的交，即 Jacobson 根，可以严格大于幂零根。第二章将详细解释这个局部化环。

\subsection{约化环与零维例子}
\label{b3:subsec:0102:03}

\(k[\varepsilon]/(\varepsilon^2)\) 只有一个素理想 \((\varepsilon)\)：素理想必须包含幂零元 \(\varepsilon\)，商为域 \(k\)。它的约化商是 \(k\)。因此只记录有哪些素理想，尚不能区分这两个环；它们的区别在于前者存在非零平方零元。

\ProseParagraphBreak
\(\mathbb Z/12\mathbb Z\) 的素理想是 \((\bar2)\)、\((\bar3)\)，其交为 \((\bar6)\)，故幂零根平方为零。约化商为 \(\mathbb Z/6\mathbb Z\cong\mathbb F_2\times\mathbb F_3\)。后者没有非零幂零元，却仍有零因子，例如 \(\bar2\bar3=0\)。约化性只排除幂零元，不排除两个不同非零元素的乘积为零。

\ProseParagraphBreak
设 \(R\) 为非零交换含幺环。若 \(R\) 的每个素理想都是极大理想，则称 \(R\) 为零维环；这等价于素理想之间不存在严格包含。一般 Krull 维数将在第九章定义。\footnote{克鲁尔（Wolfgang Krull，1899-1971），德国数学家，研究交换环中的素理想、局部化与维数。}上述双数环及有限个域的乘积都零维。

\ProseParagraphBreak
零维性却不能控制理想的降链。令 \(k\) 为域，取无限乘积环 \(R=\prod_{i\ge1}k\)，其中的理想
\[
J_n=\{\,(a_i)\in R\mid a_1=\cdots=a_n=0\,\}
\qquad(n\ge1)
\]
满足 \(J_{n+1}\subsetneq J_n\)：第 \(n+1\) 个坐标为一、其余坐标为零的元素属于 \(J_n\)，却不属于 \(J_{n+1}\)。因此 \(R\) 不是 Artin 环。

\ProseParagraphBreak
另一方面，\(R\) 的每个素理想都极大。这里不把素理想限定为坐标投影的核；无限乘积中还可以有其他素理想，下面的论证对任意素理想都适用。对任意元素 \(a=(a_i)\)，在非零坐标取 \(b_i=a_i^{-1}\)，在零坐标取 \(b_i=0\)，便有 \(a=a^2b\)。模去任一素理想后，所得商是整环；其中任意非零的 \(\bar a\) 都可从 \(\bar a=\bar a^2\bar b\) 消去一份 \(\bar a\)，得到 \(1=\bar a\bar b\)，所以这个商是域。这样，\(R\) 是零维环，却有无限严格理想降链。第三章将讨论 Artin 环及其与零维性的关系。
